\chapter{Теоретичні та математичні основи протидії злоякісному коду}
\chaptershort{Теоретичні основи протидії}
\label{chap:ch03}

Протидія злоякісному коду лише на перший погляд є суто інженерною задачею. За практичними засобами виявлення стоять фундаментальні теоретичні результати, що визначають принципові межі можливого: чи взагалі можливий досконалий детектор; якою ціною досягається задана якість виявлення; чому системи з високою точністю на практиці породжують переважно хибні тривоги; як виміряти «схожість» двох зразків і стиснутість коду. Ігнорування цих основ призводить до хибних очікувань (віра в «абсолютний антивірус»), неправильного оцінювання засобів (без урахування базової частоти) і неефективних рішень. Розділ систематизує теоретичний апарат протидії: теоретико-обчислювальні межі виявлення, теорію статистичного виявлення й робочі характеристики, проблему базової частоти, інформаційно-теоретичні міри (ентропію й подібність), моделі ознакового простору й метрики ефективності, і завершується розгорнутим числовим прикладом оцінювання детектора. Матеріал є математичним ядром посібника, на яке спираються розділи аналізу (\ref{chap:ch06}--\ref{chap:ch09}) і захисту (\ref{chap:ch10}--\ref{chap:ch11}).

\section{Постановка задачі виявлення}\label{sec:ch03-problem}

Задачу виявлення злоякісного коду формалізують як задачу класифікації. Нехай $\mathcal{X}$ --- простір об'єктів (файлів, процесів, потоків подій), кожному з яких відповідає істинний клас $y\in\{0,1\}$: $1$ --- злоякісний, $0$ --- доброякісний. \term{Детектор} --- це функція $d:\mathcal{X}\to\{0,1\}$, що наближає невідому істинну розмітку. Якість детектора визначають чотири результати: істинно-позитивні (TP, злоякісний виявлено), хибно-позитивні (FP, доброякісний помилково позначено злоякісним), істинно-негативні (TN) і хибно-негативні (FN, злоякісний пропущено). З них виводять базові показники: \term{ймовірність виявлення} (чутливість, повнота) $P_D = \mathrm{TP}/(\mathrm{TP}+\mathrm{FN})$ і \term{ймовірність хибної тривоги} $P_{FA} = \mathrm{FP}/(\mathrm{FP}+\mathrm{TN})$.

Принципово розрізняти два режими. У \emph{статичному} аналізі детектор досліджує об'єкт без його виконання (вміст файлу, структуру, рядки, ознаки); у \emph{динамічному} --- спостерігає поведінку під час виконання (виклики системних функцій, зміни в системі, мережеву активність). Кожен режим має власні теоретичні межі: статичний обмежений нерозв'язністю аналізу коду, динамічний --- неповнотою спостереження (один запуск не розкриває всіх гілок поведінки). Ці обмеження детально розглянуто в розділах~\ref{chap:ch06} і~\ref{chap:ch07}.

Порівняння теоретичних властивостей статичного й динамічного підходів наведено в табл.~\thechapter.1.

\tabhead{Таблиця~\thechapter.1 --- Теоретичні властивості статичного й динамічного виявлення}

\begin{center}\footnotesize
\begin{tabularx}{\linewidth}{L{22mm} Y Y}
\toprule
Властивість & Статичний аналіз & Динамічний аналіз \\
\midrule
Що досліджує & структуру й вміст без виконання & поведінку під час виконання \\
Повнота покриття & усі гілки коду (за ідеалу), але нерозв'язність & лише виконані гілки (неповнота) \\
Головна межа & нерозв'язність/обфускація (розділ~\ref{chap:ch08}) & антианаліз, один шлях виконання \\
Вартість & низька, швидкий & висока, потребує ізоляції \\
Ризик & відсутній (без запуску) & виконання реального коду \\
\bottomrule
\end{tabularx}
\end{center}

Слід також розрізняти \emph{виявлення} (detection) і \emph{класифікацію} (classification). Виявлення відповідає на бінарне запитання «злоякісний чи ні», тоді як класифікація визначає сімейство, тип або конкретну загрозу. Виявлення є передумовою реагування, а класифікація --- атрибуції й розуміння наміру (розділ~\ref{chap:ch16}). Обидві задачі спираються на спільний апарат, але мають різні критерії якості: для виявлення критичні повнота й частота хибних тривог, для класифікації --- матриця плутанини між класами. Ще одна важлива вісь --- \emph{момент} рішення: виявлення до виконання (превентивне, блокує запуск), під час виконання (детективне, зупиняє в процесі дії) і після (ретроспективне, за журналами й телеметрією для розслідування). Чим раніший момент, тим менша шкода, але тим менше інформації доступно й тим вища невизначеність рішення --- фундаментальний компроміс, що проходить через усі технології захисту.

\section{Теоретико-обчислювальні межі виявлення}\label{sec:ch03-limits}

\subsection{Нерозв'язність досконалого виявлення}

Фундаментальний результат теорії полягає в тому, що \emph{досконалий детектор злоякісного коду у загальному випадку неможливий}. Коен~\cite{ch01:cohen} звів задачу точного виявлення вірусів до проблеми зупинки машини Тюрінга й довів її алгоритмічну нерозв'язність. Інтуїція доведення: припустимо, існує програма $D$, що для будь-якої програми $P$ безпомилково визначає, чи є $P$ вірусом. Тоді можна побудувати «суперечливу» програму $V$, яка запитує $D$ про себе й поводиться протилежно до відповіді --- якщо $D$ каже «вірус», $V$ не розмножується, і навпаки. Це призводить до суперечності, отже $D$ не існує.

Загальнішою формою є теорема Райса~\cite{ch03:rice}: будь-яка нетривіальна семантична властивість програм (властивість поведінки, а не синтаксису) є алгоритмічно нерозв'язною. Оскільки «злоякісність» --- саме семантична властивість (що програма \emph{робить}), її точне визначення для довільної програми неможливе. Чесс і Вайт~\cite{ch03:chess} показали існування принципово невиявних вірусів, а Спінелліс~\cite{ch03:spinellis} довів, що навіть для вірусів обмеженої довжини надійна ідентифікація є NP-повною. Систематизацію нерозв'язних задач аналізу ЗПК наведено в~\cite{ch03:selcuk}.

Ці результати мають і зворотну проєкцію --- на можливості зловмисника. Та сама нерозв'язність, що обмежує захисника, обмежує й автоматичне \emph{доведення} злоякісності, але водночас дає зловмиснику теоретичне підґрунтя для ухилення: він може будувати код, аналіз якого еквівалентний нерозв'язній задачі (наприклад, через непрозорі предикати, опосередковані переходи, самомодифікацію --- розділ~\ref{chap:ch08}), змушуючи статичний аналізатор або давати невизначену відповідь, або робити консервативні припущення. Практичні статичні аналізатори обходять це апроксимаціями: аналізом у бік безпеки (over-approximation, що дає хибні тривоги) або неповним аналізом (under-approximation, що дає пропуски). Вибір напряму апроксимації --- інженерне рішення, що визначає профіль помилок інструмента. Жодна апроксимація не усуває фундаментальної межі; вона лише перерозподіляє неминучі помилки між FP і FN відповідно до призначення засобу.

\subsection{Практичні наслідки}

Нерозв'язність не означає марності виявлення --- вона означає, що будь-який реальний детектор \emph{неминуче} припускається помилок обох типів (FP і FN) і працює з наближеннями. Звідси три практичні висновки. По-перше, не існує й не може існувати «абсолютного антивірусу»; будь-яке твердження про 100\,\% виявлення є або маркетингом, або стосується скінченної тестової вибірки. По-друге, виявлення --- це завжди компроміс між $P_D$ і $P_{FA}$, і вибір робочої точки є інженерним і контекстним рішенням (підрозділ~\ref{sec:ch03-roc}). По-третє, оскільки точний аналіз неможливий, практичні засоби комбінують багато \emph{неповних} ознак (сигнатури, евристики, поведінку, статистику), кожна з яких виявляє частину загроз, а разом вони дають прийнятне покриття. Саме тому глибоко ешелонований захист (розділ~\ref{chap:ch01}) є не лише організаційним, а й математично обґрунтованим принципом.

\section{Теорія статистичного виявлення}\label{sec:ch03-detection}

\subsection{Перевірка гіпотез і критерій Неймана--Пірсона}

Виявлення зводиться до перевірки статистичних гіпотез: $H_0$ --- об'єкт доброякісний, $H_1$ --- злоякісний. Детектор обчислює за об'єктом ознакову статистику $T(x)$ і порівнює з порогом $\tau$. Згідно з лемою Неймана--Пірсона~\cite{ch03:kay} оптимальним за фіксованої $P_{FA}$ є детектор відношення правдоподібності

\begin{equation}
\Lambda(x) = \frac{p(x\,|\,H_1)}{p(x\,|\,H_0)} \ \mathop{\gtrless}_{H_0}^{H_1}\ \tau,
\label{eq:ch03-np}
\end{equation}

де $p(x\,|\,H_i)$ --- розподіли ознак за відповідної гіпотези. На практиці ці розподіли невідомі й оцінюються за навчальними даними (розділ~\ref{chap:ch09}), а $\Lambda$ наближається класифікатором. Вибір порогу $\tau$ задає компроміс: зростання $\tau$ знижує $P_{FA}$ ціною зниження $P_D$.

\subsection{Робоча характеристика приймача}\label{sec:ch03-roc}

Залежність $P_D$ від $P_{FA}$ за зміни порогу --- \term{робоча характеристика приймача} (ROC)~\cite{ch03:fawcett} --- є основним інструментом порівняння детекторів (рис.~\thechapter.1). Ідеальний детектор проходить через точку $(0,1)$; детектор, що вгадує навмання, лежить на діагоналі. Інтегральною мірою якості, незалежною від порогу, є площа під кривою (AUC). Для нормально розподілених ознак із відстанню між класами $d'$ (у стандартних відхиленнях) справедливо $P_D = Q\bigl(Q^{-1}(P_{FA}) - d'\bigr)$, де $Q$ --- доповнювальна функція стандартного нормального розподілу.

\begin{tikzfig}
\begin{tikzpicture}[font=\scriptsize,x=30mm,y=30mm]
\draw[->] (0,0) -- (1.12,0) node[below] {$P_{FA}$};
\draw[->] (0,0) -- (0,1.12) node[left] {$P_D$};
\foreach \v in {0.5,1} {\draw (\v,0) -- (\v,-0.02) node[below] {\v}; \draw (0,\v) -- (-0.02,\v) node[left] {\v};}
\draw[dashed] (0,0) -- (1,1);
\draw[thick,domain=0.001:1,samples=60] plot (\x,{1-(1-\x)^3});
\draw[thick,dotted,domain=0.001:1,samples=60] plot (\x,{1-(1-\x)^1.6});
\node[font=\tiny] at (0.55,0.93) {кращий детектор};
\node[font=\tiny] at (0.68,0.55) {гірший};
\node[font=\tiny,rotate=45] at (0.72,0.62) {випадкове вгадування};
\end{tikzpicture}
\figcap{Рис.~\thechapter.1 --- Робоча характеристика приймача (ROC) детекторів злоякісного коду}
\end{tikzfig}

\section{Проблема базової частоти}\label{sec:ch03-baserate}

Найважливіший і найчастіше ігнорований теоретичний результат для практики --- \term{проблема базової частоти}~\cite{ch03:axelsson}. Для оцінювання корисності детектора вирішальні не $P_D$ і $P_{FA}$ самі собою, а \emph{частка істинних тривог серед усіх тривог} --- позитивна прогностична цінність. За теоремою Баєса, якщо апріорна частка злоякісних об'єктів у потоці дорівнює $\pi$, то

\begin{equation}
\mathrm{PPV} = \Pr\{H_1\,|\,\text{тривога}\} = \frac{P_D\,\pi}{P_D\,\pi + P_{FA}\,(1-\pi)}.
\label{eq:ch03-bayes}
\end{equation}

Критична особливість протидії ЗПК: у реальному потоці подій злоякісні об'єкти \emph{рідкісні} ($\pi \ll 1$). Нехай система моніторингу аналізує процеси з $P_D = 0{,}99$, $P_{FA} = 0{,}001$, а злоякісним є один процес із $10^5$ ($\pi = 10^{-5}$). Тоді за~\eqref{eq:ch03-bayes} $\mathrm{PPV} = (0{,}99\cdot10^{-5})/(0{,}99\cdot10^{-5} + 10^{-3}\cdot(1-10^{-5})) \approx 0{,}0098$: понад 99\,\% тривог --- хибні. Детектор із чудовими $P_D$ і $P_{FA}$ виявляється практично непридатним через рідкісність цілі.

Наслідки для проєктування: для рідкісних загроз потрібно або різко знижувати $P_{FA}$ (на багато порядків), або підвищувати ефективну $\pi$ через попередню фільтрацію (аналізувати лише вже підозрілі події), або --- найважливіше --- \emph{корелювати кілька незалежних детекторів}. Для $k$ незалежних детекторів з частотами хибних тривог $P_{FA}^{(i)}$ правило «тривога лише за згоди всіх» дає сумарну частоту $\prod_i P_{FA}^{(i)}$, що швидко знижує знаменник~\eqref{eq:ch03-bayes}. Це математичне обґрунтування архітектури сучасних засобів (EDR, SIEM), які корелюють десятки слабких сигналів замість покладання на один «сильний» детектор (розділи~\ref{chap:ch10}, \ref{chap:ch11}).

Проблема базової частоти має й зворотний бік --- \emph{вартість хибних тривог}. Кожна хибна тривога потребує часу аналітика на розгляд; за обмежених ресурсів надмірний потік хибних тривог призводить до «втоми від тривог» (alert fatigue), коли справжні загрози губляться серед шуму або ігноруються. Це не гіпотетична, а головна операційна проблема центрів безпеки: система з високою чутливістю, але неприйнятною точністю, фактично знижує захищеність, відвертаючи увагу від реальних інцидентів. Тому проєктування детекторів для рідкісних загроз має оптимізувати не чутливість, а \emph{керованість потоку тривог} при заданому мінімальному рівні виявлення. Формально це задача оптимізації робочої точки з урахуванням пропускної спроможності аналітиків $\Lambda_{\max}$: допустима частота хибних тривог обмежена $P_{FA} \le \Lambda_{\max}/((1-\pi)\,N)$, де $N$ --- обсяг потоку. Для великих $N$ ця межа надзвичайно жорстка, що знову приводить до необхідності кореляції й автоматизації тріажу (розділ~\ref{chap:ch16}).

\section{Інформаційно-теоретичні міри}\label{sec:ch03-entropy}

\subsection{Ентропія та виявлення пакування}

Багато технік ухилення (пакування, шифрування, розділ~\ref{chap:ch08}) роблять код статистично подібним до випадкових даних. Кількісною мірою є ентропія Шеннона~\cite{ch03:shannon} байтового розподілу фрагмента $B$:

\begin{equation}
H(B) = -\sum_{i=0}^{255} p_i \log_2 p_i,\qquad 0 \le H(B) \le 8,
\label{eq:ch03-entropy}
\end{equation}

де $p_i$ --- частота байта зі значенням $i$. Звичайний виконуваний код має ентропію приблизно 6,0--6,5~біт/байт, текст --- 4--5, а стиснуті чи зашифровані дані --- близьку до 8. Лида й Гамрок~\cite{ch03:lyda} показали, що аномально висока ентропія секцій файлу є ефективним індикатором пакування чи шифрування. Проте ентропія сама собою не доводить злоякісності (легітимні програми теж пакують), тому вона є \emph{однією} з ознак у багатофакторному рішенні (розділ~\ref{chap:ch06}). Теоретичною межею є зв'язок із колмогоровською складністю~\cite{ch03:li}: нестисливість фрагмента свідчить про відсутність регулярної структури, характерної для коду.

\subsection{Міри подібності та нечітке гешування}

Оскільки ЗПК існує сімействами (варіанти одного коду), ключовою є задача вимірювання \emph{подібності} зразків. Криптографічний геш (MD5, SHA-256) непридатний: зміна одного байта повністю змінює геш. Тому застосовують \term{нечітке гешування} (context-triggered piecewise hashing, ssdeep)~\cite{ch03:kornblum}, що дає схожі геші для схожих файлів, і покускові/імпортні геші. Подібність двох об'єктів, поданих множинами ознак $A$ і $B$ (наприклад, $n$-грамів байтів~\cite{ch03:kolter} чи викликів функцій), вимірюють коефіцієнтом Жаккара:

\begin{equation}
J(A,B) = \frac{|A \cap B|}{|A \cup B|},\qquad J \in [0;1].
\label{eq:ch03-jaccard}
\end{equation}

Для ефективного обчислення подібності на великих колекціях застосовують локально-чутливе гешування (MinHash, SimHash), що зводить порівняння мільйонів зразків до пошуку близьких гешів. Міри подібності є основою кластеризації й атрибуції сімейств (розділ~\ref{chap:ch06}) і класифікації машинним навчанням (розділ~\ref{chap:ch09}).

Вибір рівня подання суттєво впливає на стійкість міри подібності до маскування. Байтові $n$-грами чутливі до пакування й перекомпіляції (поверхневі зміни руйнують подібність), тоді як подібність на рівні абстрактних структур (графа потоку керування, графа викликів, послідовностей системних викликів) стійкіша, бо відображає семантику, а не форму. Це ще один прояв «піраміди болю» (розділ~\ref{chap:ch02}): чим семантичніша ознака, тим дорожче зловмиснику її змінити, зберігши функціональність. Проте обчислення структурної подібності (наприклад, ізоморфізму підграфів) є дорогим, тому на практиці застосовують компроміси: наближені геші графів, відбитки імпортів, подібність послідовностей викликів. Задача ефективного й стійкого вимірювання подібності залишається активним напрямом досліджень, особливо в контексті великих колекцій і змагальних змін (розділ~\ref{chap:ch09}).

\section{Ознаковий простір і межі навчання}\label{sec:ch03-features}

Статистичні й машинно-навчальні детектори працюють в \term{ознаковому просторі}: об'єкт подається вектором ознак $\mathbf{x}\in\mathbb{R}^n$ (частоти $n$-грамів, імпортовані функції, ентропія секцій, послідовності викликів тощо), а детектор --- поверхнею, що розділяє класи. Теоретичні обмеження тут такі. \emph{Прокляття розмірності}: зі зростанням $n$ потрібна експоненційно більша навчальна вибірка для надійної оцінки меж. \emph{Зсув розподілу}: розподіл ознак ЗПК змінюється з часом (нові сімейства, нові техніки), тому модель, навчена на минулому, деградує --- явище \term{дрейфу концепту}, що вимагає постійного перенавчання (розділ~\ref{chap:ch09}). \emph{Змагальність}: на відміну від природних даних, зловмисник \emph{активно} змінює зразки, щоб уникнути виявлення, тому ознаки, які легко маніпулювати (поверхневі байтові), ненадійні, а стійкішими є семантичні й поведінкові (зв'язок із «пірамідою болю», розділ~\ref{chap:ch02}). Ці обмеження визначають, чому поведінкові й гібридні методи~\cite{ch03:hybrid,ch03:opirskyi-ml} переважають суто статичні для стійкого виявлення.

\section{Моделі поширення злоякісного коду}\label{sec:ch03-propagation}

Для самопоширюваного коду (черв'яків, мережевих вимагачів) теоретичну основу протидії дають \term{епідемічні моделі}. Найпростіша --- модель SI: якщо $I(t)$ --- частка заражених вузлів, $\beta$ --- інтенсивність зараження, $N$ --- загальна кількість вразливих вузлів, то на початковій стадії

\begin{equation}
\frac{dI}{dt} = \beta\, I\,\Bigl(1 - \frac{I}{N}\Bigr),
\label{eq:ch03-si}
\end{equation}

що дає логістичну криву (рис.~\thechapter.2): спершу експоненційне зростання, потім насичення. Реалістичніша модель SIR додає стан «вилікуваних» (ізольованих, пропатчених) з інтенсивністю $\gamma$, а ключовим параметром стає \term{базовий репродуктивний коефіцієнт} $R_0 = \beta/\gamma$: епідемія поширюється, якщо $R_0 > 1$, і згасає, якщо $R_0 < 1$. Ця рамка кількісно обґрунтовує заходи протидії: зменшення $\beta$ (сегментація, фільтрація, усунення вразливостей знижують імовірність успішного зараження сусіда) і збільшення $\gamma$ (швидке виявлення й ізоляція, автоматичне оновлення) переводять систему нижче порога $R_0 = 1$, зупиняючи епідемію навіть без повного охоплення захистом усіх вузлів. Історичні епідемії (Slammer, Conficker) і мережеві вимагачі (WannaCry, NotPetya, розділ~\ref{chap:ch12}) підтверджують, що швидкість поширення визначається саме $\beta$ й щільністю вразливих вузлів, а не складністю коду.

\begin{tikzfig}
\begin{tikzpicture}[font=\scriptsize,x=11mm,y=26mm]
\draw[->] (0,0) -- (4.4,0) node[below left] {час};
\draw[->] (0,0) -- (0,1.15) node[left] {$I/N$};
\draw (0,1) -- (-0.05,1) node[left] {1};
\draw[dashed] (0,1) -- (4.3,1);
\draw[thick,domain=0:4.3,samples=60] plot (\x,{1/(1+30*exp(-2.2*\x))});
\draw[thick,dotted,domain=0:4.3,samples=60] plot (\x,{1/(1+30*exp(-1.0*\x))});
\draw[thick,dashdotted,domain=0:4.3,samples=60] plot (\x,{0.033*exp(-0.25*\x)});
\node[font=\tiny,anchor=west] at (2.3,0.96) {$R_0\gg1$};
\node[font=\tiny,anchor=west] at (3.3,0.72) {$R_0>1$};
\node[font=\tiny,anchor=west] at (2.5,0.1) {$R_0<1$ (згасання)};
\end{tikzpicture}
\figcap{Рис.~\thechapter.2 --- Динаміка поширення самопоширюваного коду за різних $R_0$ (моделі SI/SIR)}
\end{tikzfig}

Докладніше моделювання поширення й ланцюгів ураження розглянуто в розділі~\ref{chap:ch04}.

\section{Метрики ефективності}\label{sec:ch03-metrics}

Для оцінювання детекторів використовують похідні від TP/FP/TN/FN метрики (табл.~\thechapter.2). Важливо вибирати метрику відповідно до задачі: за сильної незбалансованості класів (рідкісні загрози) точність (accuracy) оманлива, бо детектор, що завжди каже «доброякісний», матиме точність $1-\pi \approx 1$. Інформативнішими є повнота, точність (precision), їхнє гармонічне середнє $F_1$ і площа під кривою точність-повнота (PR-AUC), а для операційного оцінювання --- вартість помилок з урахуванням їх несиметричності.

\tabhead{Таблиця~\thechapter.2 --- Метрики ефективності виявлення злоякісного коду}

\begin{center}\footnotesize
\begin{tabularx}{\linewidth}{L{26mm} L{30mm} Y}
\toprule
Метрика & Формула & Призначення \\
\midrule
Повнота (чутливість) & $\mathrm{TP}/(\mathrm{TP}+\mathrm{FN})$ & частка виявлених загроз \\
Точність (precision) & $\mathrm{TP}/(\mathrm{TP}+\mathrm{FP})$ & частка істинних серед тривог \\
Хибно-позитивна частота & $\mathrm{FP}/(\mathrm{FP}+\mathrm{TN})$ & «шум» для аналітика \\
$F_1$-міра & $2PR/(P{+}R)$ & баланс точності й повноти \\
PPV (з базовою частотою) & рівн.~\eqref{eq:ch03-bayes} & реальна цінність тривоги \\
AUC & площа під ROC & якість без прив'язки до порогу \\
\bottomrule
\end{tabularx}
\end{center}

Окремо для протидії ЗПК критичні \emph{часові} метрики: час виявлення від моменту проникнення (визначає dwell time, розділ~\ref{chap:ch13}) і час реагування (розділ~\ref{chap:ch16}), прямо пов'язані з регуляторними вимогами щодо повідомлення (розділ~\ref{chap:ch02}).

Важливою методологічною вимогою є коректність \emph{процедури оцінювання}. Поширені помилки, що завищують якість детектора: тестування на тих самих даних, на яких навчалися (витік даних); використання збалансованих тестових вибірок, що приховує проблему базової частоти; відсутність часового розділення (навчання на пізніших зразках, тест на ранніх), що маскує дрейф концепту; відбір «легких» зразків. Коректне оцінювання вимагає часового розділення (train на минулому, test на майбутньому), реалістичної базової частоти, незалежної тестової вибірки й урахування змагальної адаптації. Для порівняння засобів галузь використовує публічні набори даних і майданчики незалежного тестування; проте результати таких тестувань слід інтерпретувати з урахуванням того, що зловмисники оптимізують свій код саме проти поширених засобів, тому висока оцінка на історичній вибірці не гарантує стійкості до нових загроз.

\section{Розгорнутий приклад оцінювання детектора}\label{sec:ch03-example}

Організація оцінює два детектори для потоку з $10^6$ файлів на добу, серед яких у середньому 50 злоякісних ($\pi = 5\cdot10^{-5}$).

\emph{Детектор A} (агресивна евристика): $P_D = 0{,}98$, $P_{FA} = 0{,}02$. Виявляє $0{,}98\cdot50 = 49$ загроз, але дає $0{,}02\cdot(10^6-50) \approx 20\,000$ хибних тривог на добу. За~\eqref{eq:ch03-bayes} $\mathrm{PPV} = 49/(49+20\,000) \approx 0{,}0024$: аналітик потоне в хибних тривогах.

\emph{Детектор B} (консервативна сигнатура): $P_D = 0{,}70$, $P_{FA} = 10^{-5}$. Виявляє 35 загроз, дає $\approx 10$ хибних тривог; $\mathrm{PPV} = 35/45 \approx 0{,}78$. Практично придатний, але пропускає 15 загроз.

\emph{Каскад} (B як перший ступінь, A як другий лише для підозрілих): якщо незалежні, правило «обидва спрацювали» дає $P_D = 0{,}98\cdot0{,}70 = 0{,}686$ і $P_{FA} = 0{,}02\cdot10^{-5} = 2\cdot10^{-7}$, звідки хибних тривог $\approx 0{,}2$/добу і $\mathrm{PPV} = 34/34{,}2 \approx 0{,}994$. Повнота при цьому впала. \emph{Паралельне правило} «будь-який спрацював» дає $P_D = 1-(1-0{,}98)(1-0{,}70) = 0{,}994$, але $P_{FA} \approx 0{,}02$ --- знову шум.

Висновок: жодне окреме правило не оптимальне; практичне рішення --- \emph{багатоступенева кореляція} з ескалацією, де перший дешевий ступінь звужує потік, а дорогі методи застосовують лише до відфільтрованого (підвищена ефективна $\pi$), плюс пріоритизація тривог. Цей приклад кількісно обґрунтовує архітектуру розділів~\ref{chap:ch10}--\ref{chap:ch11} і показує, що оцінювати детектор у відриві від базової частоти й робочого потоку некоректно.

\section*{Висновки до розділу}
\addcontentsline{toc}{section}{Висновки до розділу}

Теоретичні основи визначають принципові межі й правильні критерії протидії злоякісному коду. Фундаментальний результат --- нерозв'язність досконалого виявлення (теореми Коена й Райса): «злоякісність» є семантичною властивістю, точне визначення якої для довільної програми алгоритмічно неможливе, тому будь-який реальний детектор неминуче припускається помилок обох типів, а «абсолютного антивірусу» не існує. Теорія статистичного виявлення формалізує компроміс між виявленням і хибними тривогами через відношення правдоподібності~\eqref{eq:ch03-np} і робочу характеристику приймача. Проблема базової частоти~\eqref{eq:ch03-bayes} --- найважливіший практичний результат: для рідкісних загроз навіть детектор із чудовими показниками дає переважно хибні тривоги, що математично обґрунтовує багатоступеневу кореляцію незалежних сигналів замість покладання на один детектор. Інформаційно-теоретичні міри --- ентропія~\eqref{eq:ch03-entropy} для виявлення пакування й міри подібності~\eqref{eq:ch03-jaccard} з нечітким гешуванням для роботи із сімействами --- є основою статичного аналізу й кластеризації. Межі навчання (прокляття розмірності, дрейф концепту, змагальність) пояснюють перевагу семантичних і поведінкових ознак над поверхневими. Метрики ефективності мають вибиратися з урахуванням незбалансованості класів і реального потоку, а оцінювати детектор у відриві від базової частоти некоректно, що наочно демонструє розгорнутий приклад.

\selfcheck
\begin{enumerate}
\item Сформулюйте задачу виявлення як задачу класифікації. Визначте $P_D$ і $P_{FA}$ через TP/FP/TN/FN.
\item Чим відрізняються теоретичні межі статичного й динамічного виявлення?
\item Викладіть ідею доведення Коена про нерозв'язність виявлення вірусів. Як із цим пов'язана теорема Райса?
\item Які три практичні висновки випливають з нерозв'язності досконалого виявлення?
\item Запишіть критерій Неймана--Пірсона~\eqref{eq:ch03-np}. Як вибір порогу впливає на $P_D$ і $P_{FA}$?
\item Що таке ROC-крива (рис.~\thechapter.1) і AUC? Як за ними порівнюють детектори?
\item Сформулюйте проблему базової частоти~\eqref{eq:ch03-bayes}. Обчисліть PPV для $P_D=0{,}95$, $P_{FA}=0{,}005$, $\pi=10^{-4}$.
\item Чому кореляція незалежних детекторів розв'язує проблему базової частоти? Як це пов'язано з архітектурою EDR/SIEM?
\item Що вимірює ентропія~\eqref{eq:ch03-entropy} і чому вона є індикатором пакування? Чому висока ентропія не доводить злоякісності?
\item Чому криптографічний геш непридатний для вимірювання подібності зразків? Опишіть нечітке гешування й коефіцієнт Жаккара~\eqref{eq:ch03-jaccard}.
\item Які теоретичні обмеження має навчання в ознаковому просторі? Що таке дрейф концепту й чому ознаки ЗПК змагальні?
\item Чому точність (accuracy) оманлива для рідкісних загроз? Які метрики інформативніші (табл.~\thechapter.1)?
\item Повторіть розрахунок прикладу підрозділу~\ref{sec:ch03-example} для $\pi = 10^{-4}$ і поясніть, як змінюється придатність детекторів.
\item Дослідницьке: для набору реальних детекторів побудуйте ROC- і PR-криві на незбалансованій вибірці, оцініть AUC і проаналізуйте вплив базової частоти на операційну придатність.
\item Дослідницьке: запропонуйте схему багатоступеневої кореляції детекторів, що максимізує PPV за заданого мінімального $P_D$, і дослідіть її стійкість до дрейфу концепту та змагальних змін зразків.
\end{enumerate}
